%----------------------------------------------------------------------
%  Metadata — EDIT THIS FILE, nothing else in config/
%  thesis.tex passes these to the bossthesis class (cover page, title
%  page, declaration). scripts/build_review_docx.py reads the
%  \newcommand lines for the review copy. Open points are inline todos:
%  they sit in the title page and the declaration box, where margin notes
%  cannot go.
%----------------------------------------------------------------------
\newcommand{\thesistitle}{Understanding Privacy Harm Communication in Health Tracking Apps}
\newcommand{\thesiskind}{Bachelor's Thesis in IT Security}
\newcommand{\thesiscitetype}{Bachelor's Thesis}
\newcommand{\authorname}{Hamza Sghaier}
\newcommand{\birthdate}{9. of June 2000}
\newcommand{\birthplacename}{Mahdia, Tunisia}
\newcommand{\advisorname}{Dr.~Kavous Salehzadeh Niksirat}
\newcommand{\examinerone}{Prof.~Dr.~Carmela Troncoso}
\newcommand{\examinertwo}{\todo[inline, inlinewidth=5cm]{Second examiner.}}
\newcommand{\thechair}{Max Planck Institute for Security and Privacy (MPI-SP)}
\newcommand{\thechairaddress}{Max Planck Institute for Security and Privacy\\Universitätsstraße 140 $\cdot$ 44799 Bochum $\cdot$ Germany}
\newcommand{\startdate}{\todo[inline]{Start date set by the examination office.}}
\newcommand{\submissiondate}{\todo[inline]{Submission date.}}

% AI tools named in the declaration box. Drafted from what the repository
% records; the author must check, correct and complete every line.
\newcommand{\aitoollist}{%
\begin{itemize}
 % \item \textit{Claude (Anthropic), used through Claude Code}: writing and
 %   debugging the collection, ranking and checking scripts in
 %   \texttt{scripts/} and the generated tables; preparing the eligibility
 %   screening decisions of the app sample, which the author checks
 %   (Section~\ref{sec:results-sample}); revising and drafting passages of
 %   Chapters~4 and~5 and Appendices~C and~E.
 %   \todo[inline]{Check and complete; name the sections.}
 % \item \textit{OpenAI Codex}: \todo[inline]{The repository contains a Codex
 %   configuration; state whether and for what it was used.}
 % \item \textit{Methods digest of the literature}: \todo[inline]{Name the tool
 %   that produced the cards and summaries in 01\_Literature/AI Literature
 %   Analysis, and its use.}
 % \item \textit{Language model of the coding pipeline}: part of the method,
 %   described in Chapter~\ref{ch:methodology} and Appendix~\ref{app:prompts}.
\item \todo[inline]{AI Usage summary.}
\end{itemize}}
